\section{Probability and Markov-process framework}
\label{sec:framework}

The models in this thesis use two clocks.  Discrete generations are appropriate
when every individual is replaced at a common update, as in the Galton--Watson
process.  Continuous time is appropriate when births, deaths, absorptions and
catastrophes occur asynchronously.  Both descriptions are built from the same
three operations: condition on what happens next, exploit independence where it
is present, and package count distributions in generating functions.

\subsection{Counts, conditioning and generating functions}

All random variables are defined on a common probability space
\((\Omega,\mathcal F,\Prb)\).  For events \(B\) with \(\Prb(B)>0\), conditional
probability is
\[
  \Prb(A\mid B)=\frac{\Prb(A\cap B)}{\Prb(B)}.
\]
Repeated use of this identity, either over a partition or at the first event of a
process, gives the laws of total probability and expectation.  In particular,
first-step analysis is not a separate trick: it is ordinary conditioning applied
at the first transition.

For a non-negative integer-valued random variable \(X\), its PGF is
\begin{equation}
  G_X(z)=\E[z^X]=\sum_{k=0}^{\infty}\Prb(X=k)z^k,
  \qquad |z|\leq1.
  \label{eq:pgf-basic}
\end{equation}
The coefficients and first two moments can be recovered by differentiation:
\begin{align}
  \Prb(X=k)&=\frac{G_X^{(k)}(0)}{k!},
  &G_X(1)&=1,
  &\E[X]&=G_X'(1),
  \label{eq:pgf-identities}\\
  \Var(X)&=G_X''(1)+G_X'(1)-\bigl(G_X'(1)\bigr)^2.
  \nonumber
\end{align}
The value \(G_X(0)=\Prb(X=0)\) is especially useful because state \(0\) represents
extinction in all population models considered here.  If
\(X=X_1+\cdots+X_N\) is a sum of independent, identically distributed counts,
then
\begin{equation}
  G_X(z)=\prod_{i=1}^{N}G_{X_i}(z)=G_{X_1}(z)^N.
  \label{eq:pgf-independence}
\end{equation}
This factorisation will account for both independent founding lineages and
independent particles entering a compartment.

Joint counts require no new principle.  For
\(\boldsymbol X=(X_1,\ldots,X_d)\), define
\begin{equation}
  G_{\boldsymbol X}(z_1,\ldots,z_d)
  =\E\!\left[\prod_{j=1}^{d}z_j^{X_j}\right].
  \label{eq:multivariate-pgf}
\end{equation}
Partial derivatives at the origin recover joint state probabilities, while
derivatives at \((1,\ldots,1)\) give moments and cross-moments.  The bivariate
case will later track particles inside and outside a receiving compartment.

\subsection{Exponential clocks and continuous-time chains}

A time-homogeneous CTMC \(X=(X_t)_{t\geq0}\) on a countable state space is
specified by its generator \(Q=(q_{ij})\).  With the row convention used here,
\begin{equation}
  \Prb(X_{t+\Delta t}=j\mid X_t=i)
  =\delta_{ij}+q_{ij}\Delta t+o(\Delta t),
  \qquad \Delta t\downarrow0,
  \label{eq:generator-short-time}
\end{equation}
where \(q_{ij}\geq0\) for \(j\neq i\) and
\(q_{ii}=-\sum_{j\neq i}q_{ij}\).  Thus
\(q_i\defeq-q_{ii}\) is the total rate of leaving state \(i\).  Conditional on
the chain being in \(i\), the holding time is \(\Exp(q_i)\), after which the
next state is \(j\neq i\) with probability \(q_{ij}/q_i\).  An absorbing state
has \(q_i=0\); population zero is the main example.

This construction rests on the memoryless property of the exponential law.  If
\(T\sim\Exp(\lambda)\), then
\[
  \Prb(T>s+t\mid T>s)=\Prb(T>t)=\e^{-\lambda t}.
\]
For independent clocks \(T_r\sim\Exp(a_r)\), their minimum and its identity obey
\begin{equation}
  \min_r T_r\sim\Exp\!\left(\sum_r a_r\right),
  \qquad
  \Prb\!\left(T_j=\min_rT_r\right)=\frac{a_j}{\sum_r a_r}.
  \label{eq:exponential-race}
\end{equation}
The winning clock is independent of the minimum waiting time.  Equation
\eqref{eq:exponential-race} is the pathwise meaning of a rate model: one may draw
one holding time from the total rate and then select the event in proportion to
its rate.

Let
\(P_{ij}(t)=\Prb(X_t=j\mid X_0=i)\) and \(P(t)=(P_{ij}(t))\).  Under the usual
regularity and non-explosion conditions, the transition semigroup satisfies
\begin{equation}
  P'(t)=P(t)Q \quad\text{(forward)},
  \qquad
  P'(t)=QP(t) \quad\text{(backward)},
  \qquad P(0)=I.
  \label{eq:kolmogorov-matrix}
\end{equation}
Both equalities hold because \(P(t)=\exp(tQ)\), but their componentwise meanings
are different.  If \(\boldsymbol p(t)\) is a row vector of state probabilities,
the forward or master equation is
\begin{equation}
  \frac{\dd p_j}{\dd t}=\sum_i p_i(t)q_{ij}.
  \label{eq:master-general}
\end{equation}
It records probability entering state \(j\) through all incoming transitions and
leaving it through the diagonal rate.  This row convention is retained in every
continuous-time calculation below.

When an analytic solution is unavailable, \eqref{eq:exponential-race} also gives
an exact simulation scheme.  In state \(x\), with event propensities \(a_r(x)\),
draw a holding time from
\(\Exp(\sum_ra_r(x))\), choose event \(r\) with probability
\(a_r(x)/\sum_ka_k(x)\), update the state, and repeat.  This is the direct
Gillespie algorithm.  In this thesis simulation is used to check or evaluate an
analytical construction, not to conceal an undefined transition mechanism.

\subsection{The linear birth--death baseline}
\label{sec:linear-bd}

The linear birth--death process is the continuous-time reference model for the
chapter.  From a population of size \(n\), births occur at total rate \(n\lambda\)
and deaths at total rate \(n\mu\):
\begin{equation}
  q_{n,n+1}=n\lambda,
  \qquad q_{n,n-1}=n\mu,
  \qquad q_{nn}=-n(\lambda+\mu),
  \label{eq:linear-bd-rates}
\end{equation}
with \(0\) absorbing.  The per-capita rates are constant, so the mean follows the
deterministic equation
\begin{equation}
  \frac{\dd}{\dd t}\E[X_t]=(\lambda-\mu)\E[X_t],
  \qquad
  \E[X_t\mid X_0=N]=N\e^{(\lambda-\mu)t}.
  \label{eq:linear-bd-mean}
\end{equation}
The mean identifies the growth regime but not the probability of establishment;
that distinction is the reason for retaining the stochastic model
\cite{Allen2010}.

Write \(p_n(t)=\Prb(X_t=n\mid X_0=N)\).  From
\eqref{eq:master-general},
\begin{align}
  p_0'(t)&=\mu p_1(t),
  \label{eq:bd-master-zero}\\
  p_n'(t)&=\lambda(n-1)p_{n-1}(t)
  +\mu(n+1)p_{n+1}(t)-n(\lambda+\mu)p_n(t),
  \qquad n\geq1.
  \label{eq:bd-master}
\end{align}
Solving this infinite system term by term is unattractive.  Its time-dependent
PGF
\[
  G_N(z,t)=\sum_{n=0}^{\infty}p_n(t)z^n
\]
compresses the system into a single first-order PDE.  Multiplying
\eqref{eq:bd-master-zero}--\eqref{eq:bd-master} by \(z^n\), summing and shifting
indices gives
\begin{equation}
  \frac{\partial G_N}{\partial t}
  =(\lambda z-\mu)(z-1)\frac{\partial G_N}{\partial z},
  \qquad G_N(z,0)=z^N.
  \label{eq:bd-pgf-pde}
\end{equation}
The branching property implies \(G_N=g^N\), where \(g=G_1\).  For
\(\lambda\neq\mu\), direct solution of \eqref{eq:bd-pgf-pde} gives
\begin{equation}
  g(z,t)=
  \frac{\mu(1-z)-(\mu-\lambda z)\e^{-(\lambda-\mu)t}}
       {\lambda(1-z)-(\mu-\lambda z)\e^{-(\lambda-\mu)t}}.
  \label{eq:bd-pgf-solution}
\end{equation}
At the critical value \(\lambda=\mu\), the continuous limit of this expression is
\begin{equation}
  g(z,t)=1-\frac{1-z}{1+\lambda t(1-z)}.
  \label{eq:bd-pgf-critical}
\end{equation}
Both formulae satisfy \(g(z,0)=z\), \(g(1,t)=1\), and
\(\partial_zg(1,t)=\e^{(\lambda-\mu)t}\), which recover the initial condition,
normalisation and mean.

Setting \(z=0\) yields the finite-time extinction probability quoted later:
\begin{equation}
  \Prb(X_t=0\mid X_0=N)=
  \begin{cases}
    \displaystyle
    \left[
      \frac{\mu\bigl(1-\e^{-(\lambda-\mu)t}\bigr)}
           {\lambda-\mu\e^{-(\lambda-\mu)t}}
    \right]^N,
    &\lambda\neq\mu,\\[14pt]
    \displaystyle
    \left(\frac{\lambda t}{1+\lambda t}\right)^N,
    &\lambda=\mu.
  \end{cases}
  \label{eq:p0t}
\end{equation}
Consequently the ultimate extinction probability from one founder is \(1\) when
\(\lambda\leq\mu\) and \(\mu/\lambda\) when \(\lambda>\mu\); from \(N\)
independent founders it is the \(N\)th power of this value.

\subsection{Continuous-time quasi-stationarity}
\label{sec:cts}

The subcritical case \(\lambda<\mu\) already exhibits the distinction between
extinction and the state observed before extinction.  Put
\(\rho=\lambda/\mu<1\).  From \eqref{eq:bd-pgf-solution}, the conditional PGF
converges to
\begin{equation}
  \frac{g(z,t)-g(0,t)}{1-g(0,t)}
  \longrightarrow
  \frac{(1-\rho)z}{1-\rho z}
  \qquad (t\to\infty).
  \label{eq:ct-yaglom-pgf}
\end{equation}
The limiting conditional distribution is therefore geometric on
\(\{1,2,\ldots\}\):
\begin{equation}
  \Prb(X_\infty^{\star}=k)=(1-\rho)\rho^{k-1},
  \qquad
  \E[X_\infty^{\star}]=\frac{1}{1-\rho}
  =\frac{\mu}{\mu-\lambda}.
  \label{eq:ctsqsmean}
\end{equation}
The same limit is obtained from any fixed \(N\geq1\), because at a sufficiently
late time the event of survival is dominated by one surviving founding lineage.

For comparison with the discrete binary model, parameterise an individual event
by its probability of being a birth,
\begin{equation}
  p=\frac{\lambda}{\lambda+\mu}.
  \label{eq:rate-to-p}
\end{equation}
Then \(\rho=p/(1-p)\), and the reciprocal of the limiting conditional mean is
\begin{equation}
  A_{\mathrm c}(p)=1-\rho=\frac{1-2p}{1-p}.
  \label{eq:Acts}
\end{equation}
This continuous-time constant is explicit.  The corresponding discrete constant
is not, and the difference between them will isolate the effect of replacing
asynchronous \(\pm1\) events by synchronous binary generations.

\FloatBarrier
